% New audit repair, 17 November 2026.
% Reader/writer: GPT-6 Astra (OpenAI, ChatGPT Pro mode), effort Pro.
% This block repairs the unprinted relative-embedding step invoked at
% sphere-twists-and-the-original-smooth-six-sphere-comparison.md:707.
% It is not a quotation or a claim of a human independent review.

\subsection{The relative embedding needed for tangent-framed bordism}
\label{courseIII:pt-relative-embedding-setup}

The normal-framed Pontryagin--Thom construction is stated for embedded
bordisms.  To apply it to the group defined using abstract stable tangent
framings, the following relative construction is needed.  Every
stabilization of an existing embedding uses identically zero appended
coordinates.  The additional coordinates used in the relative extension
vanish on the prescribed boundary collars, as specified below.  All
changes of their order are specified.

\begin{lemma}[Relative Euclidean embedding with specified product collars]
\label{courseIII:pt-relative-embedding}
Let $W$ be a compact smooth manifold whose boundary is the disjoint union
$M_0\sqcup M_1$.  Fix disjoint collars and smooth embeddings
$j_i:M_i\longrightarrow\mathbb R^D$.  After appending finitely many zero
spatial coordinates, there is a smooth embedding
\[
 J:W\longrightarrow\mathbb R^N\times[0,1]
\]
whose restrictions in smaller collars are exactly
$(m,t)\mapsto(j_0(m),0,t)$ and
$(m,t)\mapsto(j_1(m),0,t)$, respectively.  Its interior is in
$\mathbb R^N\times(0,1)$.  Here the coordinate in the upper collar is
$t=1-r$, where $r$ is the prescribed inward collar coordinate.
\end{lemma}

\begin{proof}
Choose a number $a>0$ so that the two collars are defined for $0\le r<4a$
and $4a<1/2$.  Choose a smooth function $t:W\to[0,1]$ equal to $r$ in
the lower collar for $r<3a$, equal to $1-r$ in the upper collar for
$r<3a$, and taking values strictly between zero and one elsewhere.
For example, interpolate by a smooth cutoff between each indicated
collar formula and the constant $1/2$, completing the interpolation
before the end of that collar.  Convex combinations remain in $(0,1)$
at every interior point.

The boundary maps have a smooth extension $f:W\to\mathbb R^D$ which is
exactly $j_i(m)$ for $r<3a$: multiply each $j_i$ in its collar by a
cutoff equal to one there and zero near the far end of the collar,
extend by zero, and add the two terms.  Their supports are disjoint.

There is a smooth embedding $e:W\to\mathbb R^K$ given by finite
coordinate blocks $(\chi_\ell,\chi_\ell\psi_\ell)$, with each cutoff
supported in a coordinate chart and identically one on a smaller cover.
Boundary charts are charts in a half-space and their coordinate
functions and cutoffs extend smoothly to open Euclidean charts.  A
block with $\chi_\ell=1$ separates points by $\psi_\ell$ and has
injective tangent derivative.  The resulting map is an injective
immersion, and compactness makes it an embedding.  This is the same
finite-chart proof as the closed-manifold construction, with the
half-space charts retained.

Choose a smooth function $g\geq0$ which is zero exactly on the two
closed collar regions $r\leq a$ and positive elsewhere.  Such a
function is obtained by a smooth flat step in each collar, equal to
one before its far end, and by setting it equal to one off those
collars.  Define
\[
 J(w)=\bigl(f(w),g(w),g(w)e(w),t(w)\bigr).
 \tag{RB.1}
\]
At a point with $g>0$, the first scalar $g$ and the block $ge$ recover
$e$ by division.  Equality of two images therefore implies equality
of the two points.  Also $dJ(v)=0$ gives $dg(v)=0$, then $g\,de(v)=0$,
so $v=0$.  Points with $g=0$ cannot have the same image as points with
$g>0$.  On $g=0$, formula (RB.1) is the prescribed product embedding;
its two collar components have disjoint height ranges.  Its derivative
is injective there because $f=j_i$ and $t=r$ or $1-r$ on a larger
collar.  These observations prove injectivity and immersion everywhere,
including the flat joining locus $r=a$.  Compactness again gives an
embedding.  The previously chosen height function proves the assertion
about the interior.  One further zero spatial coordinate may be
appended if an even value of $N$ is wanted.
\end{proof}

\begin{lemma}[Exact boundary normal frame of a tangent-framed bordism]
\label{courseIII:pt-relative-normal-frame}
Let $\dim W=7$.  Suppose the stable tangent framing of $W$ is a product
in the specified collars and induces the prescribed stable frames
\[
 \tau_i:TM_i\oplus\mathbb R^k\longrightarrow\mathbb R^{6+k}.
\]
Append an identity line if necessary so that $k$ is even.  In a collar
use the coordinate vector $\partial_t$ directed from the lower end
towards the upper end, and write the product framing as
\[
 \tau_W(v+a\partial_t,b)=(\tau_i(v,b),a).
 \tag{RB.2}
\]
The embedding in the preceding lemma, after a further specified zero
stabilization, has a normal framing whose restriction is exactly the
normal frame constructed from $j_i$ and $\tau_i$, followed by one
identity normal vector.  Hence it is an embedded normal-framed bordism
between those stabilized endpoint frames.
\end{lemma}

\begin{proof}
Regard the derivative of $J$ as a map
$K:TW\to\mathbb R^N\oplus\mathbb R_t$, put
$Q_W=(K^{\mathsf T}K)^{-1}K^{\mathsf T}$ in any tangent chart, and put
$P_W=I-KQ_W$.  The matrix $P_W$ is independent of the tangent chart
and is the orthogonal projection onto the normal bundle $\nu_W$.
The identities $Q_WK=I$, $Q_WP_W=0$, $P_W+KQ_W=I$ prove all of these
assertions directly.

Add $7+k$ zero ambient coordinates.  The normal bundle is then
$\nu_W\oplus\mathbb R^{7+k}$.  Its following trivialization and
inverse are actual bundle maps:
\[
 \begin{aligned}
 A_W(y,u)&=(y+Kv,b), &&(v,b)=\tau_W^{-1}(u),\\
 A_W^{-1}(z,b)&=(P_Wz,\tau_W(Q_Wz,b)).
 \end{aligned}
 \tag{RB.3}
\]
Substitution using the three projection identities proves both inverse
identities.  This also proves that the formulas agree in overlapping
tangent charts.  The images under $A_W^{-1}$ of a fixed ordered basis
give a smooth full normal frame, with its actual Gram matrix; no
orthonormal replacement is needed.

Put the new ambient coordinates before the height coordinate, so the
ambient cylinder has spatial order
$\mathbb R^N,\mathbb R^{6+k},\mathbb R_a$ and last coordinate $t$.
Order the normal parameters as $(z,b,a)$, with respective dimensions
$N,k,1$.  This parameter order is the explicitly recorded permutation
of the order $(z,a,b)$ in the right side of (RB.3).  In a product collar
write $J_i=Dj_i$, $Q_i=(J_i^{\mathsf T}J_i)^{-1}J_i^{\mathsf T}$,
and $P_i=I-J_iQ_i$.  The derivative and projection of the original
bordism embedding are block diagonal there.  Applying (RB.2)--(RB.3)
gives the full collar frame
\[
 (z,b,a)\longmapsto
 \bigl(P_i z,\tau_i(Q_i z,b),a;0\bigr).
 \tag{RB.4}
\]
The final semicolon separates the height coordinate.  The first two
blocks in (RB.4) are precisely the endpoint frame in the original
tangent-to-normal construction.  The third is its right-appended unit
vector.  The last component is zero throughout the collar, so the
embedding and the normal frame have the exact product germs required
by the embedded Pontryagin--Thom construction.  Formula (RB.4), rather
than an orientation-only argument, verifies the asserted matching.

If the initial definition of tangent-framed bordism supplies a framing
homotopy on the boundary instead of the exact product equality
(RB.2), make that homotopy constant near both endpoints and install it
on a smaller collar.  It is a homotopy through bundle isomorphisms,
equals the old frame on a neighbourhood of the inner seam, and hence
pastes to the old framing on the remaining manifold.  This gives
(RB.2) exactly and does not change the specified boundary maps.
\end{proof}

\begin{lemma}[The inverse normal-frame comparison]
\label{courseIII:pt-normal-inverse-comparison}
Starting with a framed normal bundle, passing to its tangent-plus-normal
stable tangent framing, and applying the preceding normal construction
recovers the original normal framing up to stabilization and an explicit
homotopy through full normal frames.
\end{lemma}

\begin{proof}
Let $j:M^6\to\mathbb R^D$ have derivative $J$ and normal frame
$E:\mathbb R^r\to\nu$, where $D=6+r$.  Suspend once if necessary
so that $r$ and $D$ are even.  Put
\[
 Q=(J^{\mathsf T}J)^{-1}J^{\mathsf T},\qquad
 P=I-JQ,
 \qquad E^+=(E^{\mathsf T}E)^{-1}E^{\mathsf T}.
\]
Then $EE^+=P$, and the stable tangent frame is
$\tau(v,b)=Jv+Eb$.  Applying its normal construction to $(j,0)$ gives
\[
 F(a,b)=(Pa,JQa+Eb)
       \quad\hbox{in }\nu\oplus\mathbb R^D.
 \tag{RB.5}
\]
The original stabilized normal frame is
$E_{\rm st}(b,a)=(Eb,a)$.  Decompose $a=Jv+Ec$ using
$v=Qa$, $c=E^+a$.  On the fixed even-dimensional parameter space
$\mathbb R^r$ choose the constant operator $I_r$ with square $-I$,
acting on successive coordinate pairs by $(x,y)\mapsto(-y,x)$.
For $0\le s\le1$, set $\theta=\pi s/2$, $\varphi=\pi s$, and
\[
 \begin{aligned}
 b_s&=\cos\theta\,b+\sin\theta\,c,\\
 c_s&=(\cos\varphi\,I+\sin\varphi\,I_r)
             (-\sin\theta\,b+\cos\theta\,c),\\
 G_s(b,a)&=(Eb_s,Jv+Ec_s).
 \end{aligned}
 \tag{RB.6}
\]
Each factor in this formula is invertible: use the inverse plane
rotation on $(b,c)$ and replace the sine by its negative in the factor
involving $I_r$.  Thus $G_s$ is a full frame for every $s$.  At $s=0$
it is $E_{\rm st}$; at $s=1$ it is $(Ec,Jv+Eb)=F(a,b)$.
The domain order $(b,a)$ at this endpoint is the recorded block
permutation of the order $(a,b)$ in (RB.5).  Its sign is
$(-1)^{Dr}=+1$.  It has a path through constant orientation-preserving
matrices to the identity: use plane rotations to reduce its orthogonal
matrix to a diagonal matrix, then rotate each pair of remaining
negative diagonal entries through angle $\pi$.  Reversing the finite
list gives the required path.  This supplies the parameter-order
homotopy as well as the bundle-dependent homotopy (RB.6).
All occurrences of $J,E,Q,E^+$ remain the given point-dependent
matrices.  Smoothness follows from their displayed formulas and matrix
inversion on positive definite Gram matrices.
\end{proof}

\begin{proposition}[The abstract tangent-framed Pontryagin--Thom comparison]
\label{courseIII:pt-abstract-framed-comparison}
The embedded normal-framed Pontryagin--Thom bijection, together with the
displayed stable tangent-to-normal construction, induces an isomorphism
$\Omega^{\rm fr}_6\cong\pi_6(\mathbb S)$ for abstract stable
tangent-framed bordism.  It is independent of the embedding and the
specified zero stabilizations, and preserves disjoint-union addition.
\end{proposition}

\begin{proof}
\paragraph{Choice of the original embedding.}
First consider appending $d$ zero coordinates to an embedding
$j:M\to\mathbb R^D$.  Its canonical frame, in the displayed orders,
is
\[
 ((a,c),b)\longmapsto(Pa,c,\tau(Qa,b)).
\]
Move the middle ambient block $c$ after the $\tau$ block and move its
parameter block after $b$.  The result is exactly the original frame
followed by $d$ identity vectors.  The two permutation signs are
$(-1)^{d(6+k)}$ and $(-1)^{dk}$; with the stipulated even $k$ they
are both positive.  The constant rotation construction in the preceding
proof supplies both homotopies.  Thus a zero-coordinate enlargement of
the original embedding agrees with the specified stable transition map,
with the complete block comparison retained.

The same comparisons account for embedding choices.  For two embeddings
$j_0,j_1$ of one compact manifold, first place them in a common Euclidean
space by zero coordinates.  In its double use
$j_s=(\cos\theta(s)\,j_0,\sin\theta(s)\,j_1)$, where $\theta$
goes from zero to $\pi/2$ and is constant near the endpoints.  At
every parameter at least one coefficient is nonzero, so injectivity
and injectivity of the tangent derivative follow from that embedding.
Apply the actual Gram-matrix normal construction, including its specified
zero stabilization, to this family; write the resulting stabilized
embedding again as $j_s$.  It gives a smooth family of full normal
frames $E_s$.  A family $E_s$ orthogonal to $Dj_s$ yields the
normal frame
\[
 E_s A\longmapsto
       (E_s A,-\langle E_s A,\partial_sj_s\rangle)
\]
on the embedded cylinder $(m,s)\mapsto(j_s(m),s)$: its scalar product
with every vector $(Dj_s v+a\partial_sj_s,a)$ is zero.  Its rank is
unchanged because its spatial part is injective.  It is a product frame
near the ends.  An ambient plane rotation carries the second coordinate
block to the first, with its full normal derivative.  Naturality of the
canonical frame under a constant orthogonal transformation $R$ reads
\[
 E_{Rj}(a,b)=(R\oplus I)E_j(R^{-1}a,b).
\]
This follows by substituting $J\mapsto RJ$, $P\mapsto RPR^{-1}$
and $Q\mapsto QR^{-1}$ in its complete formula.  The plane rotation
path consequently supplies the parameter comparison as well as the
ambient one.  The positive block-order comparisons are the
constant-matrix homotopies already recorded above.  This gives embedded framed bordisms for all
embedding and ordering choices, after finite stabilization.

The group operation is also preserved.  Place two disjoint compact
embedded representatives in disjoint Euclidean balls by translations
and positive dilations, retaining their full transported frames.  Their
tubes can be chosen inside those balls.  The collapse of their disjoint
union is the map supported on those two balls, constant elsewhere.
The based pinch description of addition in the sphere homotopy group
is exactly this map.  The translations and dilations give framed
cylinders by the same family construction, so this comparison is
independent of the placements.

Consequently the two existing Pontryagin--Thom constructions descend to
abstract tangent-framed bordism in dimension six.  The tangent composite
is the explicit involution homotopy already proved in the sphere-twist
chapter, equations (7.9)--(7.19).  The normal composite is
(RB.5)--(RB.6).  Relative embedding and (RB.4) give bordism compatibility.
Together with the already proved embedded normal-framed bijection, these
prove the asserted isomorphism
$\Omega^{\rm fr}_6\cong\pi_6(\mathbb S)$ at the precise point where
the geometric framing chapter uses it.  This bridge does not compute
the stable group and does not certify any still unread Adams prerequisite.
\end{proof}
